\documentclass[11pt]{article}

\usepackage[margin=1in]{geometry}
\usepackage[T1]{fontenc}
\usepackage{lmodern}
\usepackage{microtype}
\usepackage{amsmath,amssymb,mathtools,bm}
\usepackage{booktabs,tabularx,array}
\usepackage{xcolor}
\usepackage{float}
\usepackage{hyperref}
\usepackage{enumitem}
\usepackage{fancyhdr}
\usepackage{listings}
\usepackage{titlesec}

\definecolor{accent}{RGB}{25,75,115}
\definecolor{lightaccent}{RGB}{235,244,249}
\definecolor{codebg}{RGB}{247,248,249}
\definecolor{warning}{RGB}{145,80,20}

\hypersetup{
  colorlinks=true,
  linkcolor=accent,
  urlcolor=accent
}

\titleformat{\section}{\Large\bfseries\color{accent}}{\thesection}{0.7em}{}
\titleformat{\subsection}{\large\bfseries\color{accent}}{\thesubsection}{0.7em}{}

\setlength{\parindent}{0pt}
\setlength{\parskip}{5pt}
\setlength{\headheight}{14pt}
\pagestyle{fancy}
\fancyhf{}
\fancyhead[L]{\footnotesize Numerical solution of the \(b\)-space jet RG}
\fancyhead[R]{\footnotesize \thepage}

\lstset{
  basicstyle=\ttfamily\small,
  backgroundcolor=\color{codebg},
  frame=single,
  columns=fullflexible,
  keepspaces=true,
  showstringspaces=false,
  breaklines=true,
  aboveskip=9pt,
  belowskip=8pt
}

\newcommand{\dd}{\mathrm{d}}
\newcommand{\ee}{\mathrm{e}}
\newcommand{\as}{a_s}
\newcommand{\Jvec}{\bm J}
\newcommand{\Pmat}{\bm P}
\newcommand{\Rvec}{\bm{\mathcal R}}
\newcommand{\bzero}{b_0}

\title{\textbf{Numerical Solution of the \texorpdfstring{\(b\)}{b}-Space Jet RG}\\
\large Problem formulation, transformed geometry, and stepwise method}
\author{}
\date{\today}

\begin{document}

\maketitle

\begin{center}
\fcolorbox{accent}{lightaccent}{
\begin{minipage}{0.92\textwidth}
\textbf{Objective.}
Given quark and gluon jet functions on the \(b\)-dependent canonical line
\(\mu=\mu_*(b)\), construct the solution throughout the finite
\((b,\mu)\) rectangle. The implementation uses a uniform grid in
\(\ell=\ln b\), derives its time grid from the same boundary line, precomputes
the complete splitting convolution, evolves both sides of the line, and
returns a callable interpolator for \(J_q(b,\mu)\) and \(J_g(b,\mu)\).
\end{minipage}}
\end{center}

\tableofcontents

\section{The \texorpdfstring{\(b\)}{b}-space evolution problem}

\subsection{Renormalization-group equation}

The object to be evolved is the two-component row vector
\begin{equation}
  \Jvec(b,\mu)
  =
  \bigl(J_q(b,\mu),J_g(b,\mu)\bigr).
  \label{eq:jet-vector}
\end{equation}
Its position-space renormalization-group equation is
\begin{equation}
  \frac{\partial}{\partial\ln\mu^2}\Jvec(b,\mu)
  =
  \int_0^1 \dd y\,y^2\,
  \Jvec\left(\frac{b}{y},\mu\right)\Pmat(y,\mu),
  \label{eq:bspace-rg}
\end{equation}
where the time-like splitting matrix is
\begin{equation}
  \Pmat(y,\mu)
  =
  \begin{pmatrix}
    P_{qq}(y,\mu) & P_{qg}(y,\mu)\\
    P_{gq}(y,\mu) & P_{gg}(y,\mu)
  \end{pmatrix}.
  \label{eq:splitting-matrix}
\end{equation}
The row-vector convention fixes the channel orientation. In components,
\begin{align}
  \frac{\partial J_q}{\partial\ln\mu^2}
  &=
  \int_0^1\dd y\,y^2
  \left[
    J_q(b/y,\mu)P_{qq}(y,\mu)
    +J_g(b/y,\mu)P_{gq}(y,\mu)
  \right],
  \label{eq:bspace-q}
  \\
  \frac{\partial J_g}{\partial\ln\mu^2}
  &=
  \int_0^1\dd y\,y^2
  \left[
    J_q(b/y,\mu)P_{qg}(y,\mu)
    +J_g(b/y,\mu)P_{gg}(y,\mu)
  \right].
  \label{eq:bspace-g}
\end{align}

\subsection{Boundary data on a scale curve}

The initial condition is not specified at one common value of \(\mu\).
Instead, every \(b\) has a canonical starting scale
\begin{equation}
  \mu_*(b)=\frac{\bzero}{b_*(b)},
  \qquad
  \bzero=2\ee^{-\gamma_E},
  \label{eq:canonical-scale}
\end{equation}
and the input is
\begin{equation}
  \Jvec\bigl(b,\mu_*(b)\bigr)
  =
  \Jvec_{\mathrm{bc}}(b).
  \label{eq:boundary-value}
\end{equation}
The numerical solver treats both \(b_*(b)\) and
\(\Jvec_{\mathrm{bc}}(b)\) as user-supplied functions. A common choice is
\begin{equation}
  b_*(b)
  =
  \frac{b}{\sqrt{1+b^2/b_{\mathrm{sat}}^2}},
  \label{eq:bstar-example}
\end{equation}
where \(b_{\mathrm{sat}}\) is the saturation parameter. This notation keeps
it distinct from the numerical endpoint \(b_{\mathrm{grid,max}}\).

The domain ordering below assumes that \(\mu_*(b)\) strictly decreases as
\(b\) increases, or equivalently that \(t_*(\ell)\) strictly decreases with
\(\ell\). The lattice constructor checks the corresponding discrete
monotonicity condition.

For input bounds
\(b_{\mathrm{grid,min}}\le b\le b_{\mathrm{grid,max}}\), the desired output
is the full solution over
\begin{equation}
  \mu_{\min}=\mu_*(b_{\mathrm{grid,max}})
  \le \mu \le
  \mu_{\max}=\mu_*(b_{\mathrm{grid,min}}).
  \label{eq:mu-domain}
\end{equation}
Thus the scale range is determined by the \(b\) range and \(b_*(b)\); it is
not an independent solver input.

Numerical examples use \(b\) and \(b_*\) in \(\mathrm{GeV}^{-1}\) and \(\mu\)
in \(\mathrm{GeV}\). Logarithms of dimensional quantities are shorthand for
ratios to the corresponding reference unit.

\section{Change of variables}

Define
\begin{equation}
  \ell=\ln b,
  \qquad
  t=\ln\mu^2,
  \qquad
  y=\ee^{-u},
  \quad u\ge0.
  \label{eq:variables}
\end{equation}
Then
\[
  \frac{b}{y}=b\ee^u,
  \qquad
  \ln\left(\frac{b}{y}\right)=\ell+u,
  \qquad
  y^2\dd y=-\ee^{-3u}\dd u.
\]
Reversing the integration limits gives
\begin{equation}
  \partial_t\Jvec(\ell,t)
  =
  \int_0^\infty\dd u\,\ee^{-3u}\,
  \Jvec(\ell+u,t)\Pmat(\ee^{-u},t).
  \label{eq:transformed-rg}
\end{equation}
The important structural fact is that the derivative at \(\ell\) depends
only on the same point and on points at larger \(\ell\), or equivalently at
larger \(b\).

In the transformed plane, the boundary curve is
\begin{equation}
  t_*(\ell)
  =
  2\ln\left[\frac{\bzero}{b_*(\ee^\ell)}\right].
  \label{eq:general-boundary-curve}
\end{equation}
For the example in Eq.~\eqref{eq:bstar-example},
\begin{equation}
  t_*(\ell)
  =
  2\ln\bzero-2\ell
  +\ln\left(1+\frac{\ee^{2\ell}}{b_{\mathrm{sat}}^2}\right).
  \label{eq:example-boundary-curve}
\end{equation}

\pagebreak
\section{The transformed numerical problem}

\subsection{A curved initial-value problem}

The input values lie on \(t=t_*(\ell)\), while the requested solution fills a
rectangle. This divides the rectangle into two regions:
\begin{align}
  t\ge t_*(\ell)
  &\quad\text{(upper region, forward evolution)},\\
  t\le t_*(\ell)
  &\quad\text{(lower region, backward evolution)}.
\end{align}
Solving only by activation and forward evolution fills the upper region but
leaves the lower region empty. A separate backward pass is therefore needed
to obtain an interpolator over the full plane.

\subsection{Nonlocality is one-sided}

Equation~\eqref{eq:transformed-rg} samples \(\ell+u\) with \(u\ge0\).
If the \(b_*\) prescription makes \(t_*(\ell)\) decrease with \(\ell\), the
larger-\(\ell\) values needed by a newly activated node have already been
activated and evolved. This one-sided dependence gives the discrete system a
triangular structure and makes explicit marching possible.

\subsection{Endpoint distributions}

The diagonal splitting functions are distributions, not ordinary functions.
The implementation uses the decomposition
\begin{equation}
  P_{aa}(y,t)
  =
  P_{aa}^{\mathrm{reg}}(y,t)
  +P_{aa}^{\delta}(t)\delta(1-y)
  +D_{aa}^{(0)}(y,t)\left[\frac{1}{1-y}\right]_+
  +D_{aa}^{(1)}(y,t)
  \left[\frac{\ln(1-y)}{1-y}\right]_+,
  \label{eq:kernel-decomposition}
\end{equation}
for \(a=q,g\). The delta terms and the self-subtraction implied by the plus
prescription must be handled separately from the sampled \(0<y<1\) values.
Treating Eq.~\eqref{eq:kernel-decomposition} as an ordinary pointwise kernel
would give the wrong evolution.

\subsection{Finite-domain closure}

The exact equation can request arbitrarily large \(b/y\). On a finite lattice,
the current solver imposes
\begin{equation}
  \Jvec(b,\mu)=0
  \qquad
  \text{for }b>b_{\mathrm{grid,max}}.
  \label{eq:zero-closure}
\end{equation}
This zero closure is part of the numerical approximation. The user must
choose \(b_{\mathrm{grid,max}}\) far enough into the region where the physical
jet function is negligible. A convenient \(b_*\) saturation scale is not, by
itself, a valid reason to truncate the jet function at the same value.

There is an additional endpoint condition from the diagonal plus
distributions. If \(\Jvec(b_{\mathrm{grid,max}},t)\) is nonzero while all
larger-\(b\) values are set abruptly to zero, the shifted real contribution
vanishes for \(y<1\) but the virtual self-subtraction at \(y=1\) remains. The
result is logarithmically sensitive to the first \(u\)-grid spacing. A
grid-convergent use of the sharp closure therefore requires the jet function
to be numerically zero at the endpoint, not merely outside it. The code does
not enforce this condition; it must be checked by the user.

\section{General numerical idea}

The method combines a lattice-aligned convolution with explicit evolution on
both sides of the boundary curve.

\subsection{One lattice for both variables}

Choose \(N\ge2\) points uniformly in \(\ell\), ordered from large to small
\(b\):
\begin{align}
  \ell_i
  &=
  \ell_{\max}-(i-1)\Delta\ell,
  &
  \Delta\ell
  &=
  \frac{\ell_{\max}-\ell_{\min}}{N-1},
  \label{eq:ell-grid}\\
  b_i&=\ee^{\ell_i},
  &
  t_i&=t_*(\ell_i),
  \qquad i=1,\ldots,N.
  \label{eq:paired-grid}
\end{align}
The time grid is exactly \((t_1,\ldots,t_N)\). There is one time node for each
\(\ell\) node, with no sorting, duplicate removal, extra regular time grid, or
independent final scale. The code requires \(t_{i+1}>t_i\).

Sample the convolution variable with the same spacing,
\begin{equation}
  u_m=m\Delta\ell,
  \qquad
  y_m=\ee^{-m\Delta\ell}.
  \label{eq:u-grid}
\end{equation}
For \(m=1,\ldots,i-1\),
\begin{equation}
  \ell_i+u_m=\ell_{i-m}.
  \label{eq:lattice-alignment}
\end{equation}
Every shifted argument therefore lands exactly on a stored \(b\) node. No
interpolation is performed inside the convolution.

\subsection{Perturbative kernel tables}

The splitting matrix is assembled as a perturbative series in
\begin{equation}
  \as(t)=\frac{\alpha_s(\ee^{t/2})}{4\pi}:
\end{equation}
\begin{equation}
  \Pmat(y,t)
  =
  \as\Pmat^{(0)}(y)
  +\as^2\Pmat^{(1)}(y)
  +\as^3\Pmat^{(2)}(y)+\cdots.
  \label{eq:perturbative-series}
\end{equation}
The public argument \texttt{order} truncates this series:
\begin{equation}
\begin{array}{c|l}
\texttt{order} & \text{included terms}\\
\hline
0 & \as\Pmat^{(0)}\\
1 & \as\Pmat^{(0)}+\as^2\Pmat^{(1)}\\
2 & \as\Pmat^{(0)}+\as^2\Pmat^{(1)}+\as^3\Pmat^{(2)}
\end{array}
\label{eq:order-map}
\end{equation}
At every \(t_n\), the code evaluates
\(\alpha_s(\mu_n)\) using \texttt{alpha\_s\_func(mu\_n)} and precomputes
all channel kernels at the non-endpoint samples \(y_m<1\). Delta coefficients
and the \(y=1\) values needed by the plus subtraction are stored separately.
Midpoint RK2 additionally precomputes a table at \((t_n+t_{n+1})/2\).

\subsection{Discrete right-hand side}

After the regular, distributional, and quadrature factors have been folded
into the kernel tables, the node-level right-hand side has the compact form
\begin{align}
  \mathcal R_{q,i}^{(n)}
  &=
  J_{q,i}^{(n)}
  \left(\Delta_{qq}^{(n)}-S_{qq}^{(n)}\right)
  +
  \sum_{m=1}^{i-1}
  \left[
    J_{q,i-m}^{(n)}K_{qq,m}^{(n)}
    +J_{g,i-m}^{(n)}K_{gq,m}^{(n)}
  \right],
  \label{eq:discrete-rhs-q}\\
  \mathcal R_{g,i}^{(n)}
  &=
  J_{g,i}^{(n)}
  \left(\Delta_{gg}^{(n)}-S_{gg}^{(n)}\right)
  +
  \sum_{m=1}^{i-1}
  \left[
    J_{q,i-m}^{(n)}K_{qg,m}^{(n)}
    +J_{g,i-m}^{(n)}K_{gg,m}^{(n)}
  \right].
  \label{eq:discrete-rhs-g}
\end{align}
Here \(\Delta_{aa}^{(n)}=P_{aa}^{\delta}(t_n)\). For an off-diagonal channel,
\(K_{ab,m}^{(n)}=\Delta\ell\,y_m^3P_{ab}(y_m,t_n)\). For a diagonal channel,
the same table entry contains
\[
  \Delta\ell\,y_m^3
  \left[
    P_{aa}^{\mathrm{reg}}(y_m,t_n)
    +\frac{D_{aa}^{(0)}(y_m,t_n)}{1-y_m}
    +\frac{D_{aa}^{(1)}(y_m,t_n)\ln(1-y_m)}{1-y_m}
  \right].
\]
The quantity \(S_{aa}^{(n)}\) is the matching endpoint self-subtraction
constructed from \(D_{aa}^{(0,1)}(1,t_n)\), including the analytically
integrated subtraction tail beyond the last sampled shift. Thus the
convolution routine only takes jet vectors and returns scalar dot products
for each output component.

\subsection{Upper-region evolution}

At time \(t_n\), nodes \(1,\ldots,n\) are active. Node \(n\) is initialized
with its exact boundary value, and all active nodes are advanced from \(t_n\)
to \(t_{n+1}\). Forward Euler uses
\begin{equation}
  \Jvec^{(n+1)}
  =
  \Jvec^{(n)}+\Delta t_n\Rvec(t_n,\Jvec^{(n)}),
  \label{eq:euler}
\end{equation}
while midpoint RK2 uses
\begin{align}
  \bm k_1&=\Rvec(t_n,\Jvec^{(n)}),\\
  \bm k_2&=\Rvec\left(
    t_n+\frac{\Delta t_n}{2},
    \Jvec^{(n)}+\frac{\Delta t_n}{2}\bm k_1
  \right),\\
  \Jvec^{(n+1)}
  &=\Jvec^{(n)}+\Delta t_n\bm k_2.
  \label{eq:rk2}
\end{align}
This pass fills all entries on and above the boundary diagonal.

The current activation routine uses an absolute time tolerance of
\(10^{-12}\). The exact \(1,\ldots,n\) statement therefore assumes adjacent
time nodes are separated by more than this tolerance, as they are in the
tested grids.

\subsection{Lower-region evolution}

The lower region is filled after the forward pass. Starting from the complete
top-time column \(t_N\), the algorithm visits columns in the order
\[
  t_N,\ t_{N-1},\ \ldots,\ t_1.
\]
At step \(n\), the values in column \(t_n\) are already known: nodes
\(i\le n\) came from the upper pass and nodes \(i>n\) came from earlier
backward steps. Nodes \(i\ge n\) are then evolved through the negative time
increment \(t_{n-1}-t_n\). This fills the next lower column without changing
the stored upper-region values or the diagonal boundary data.

The descending-column construction is possible because
Eqs.~\eqref{eq:discrete-rhs-q}--\eqref{eq:discrete-rhs-g} only refer to node
indices smaller than or equal to \(i\). Euler and midpoint RK2 are implemented
with the same sign convention as in the forward pass.

\subsection{Boundary-preserving interpolation}

The stored solution is an \(N\times N\) history for each channel. Away from a
cell crossed by the boundary diagonal, ordinary bilinear interpolation is
used. A diagonal cell is split into an upper and a lower triangle, and the
interpolator uses only the three vertices from the relevant side. This avoids
mixing independently evolved upper and lower values.

Consequently, every lattice value is returned exactly, and the boundary is
preserved exactly on the piecewise-linear diagonal connecting its lattice
nodes. Between nodes, this is the lattice representation of the continuous
curve, not an additional exact evaluation of the user boundary function.

\section{Numerical implementation and use}
\label{sec:numerical-implementation}

\subsection{Code organization}

The implementation is divided into small files with one primary
responsibility each.

\begin{table}[H]
\centering
\caption{Numerical RG source files.}
\label{tab:source-files}
\begin{tabularx}{\textwidth}{@{}>{\ttfamily}l X@{}}
\toprule
\textrm{File} & \textrm{Responsibility}\\
\midrule
numerical\_rg.jl
  & Loads the constants, running coupling, splitting functions, and numerical
    solver modules. This is the entry file for users.\\
grids.jl
  & Builds the paired \(\ell_i\), \(b_i\), \(b_{*,i}\), \(\mu_i\), and \(t_i\)
    arrays in the \texttt{LatticeGrid} structure, and samples boundary data.\\
components.jl
  & Builds \texttt{SplittingKernelTable} objects at node or midpoint times,
    including regular, delta, and plus-distribution contributions.\\
convolution.jl
  & Evaluates the quark and gluon right-hand sides from precomputed kernel
    rows and jet vectors. Each convolution is a serial SIMD dot product.\\
solver.jl
  & Performs activation, upper-region evolution, lower-region
    descending-column evolution, and constructs the solution object.\\
interpolation.jl
  & Defines \texttt{StepwiseRGSolution}, domain checks, boundary-scale
    interpolation, and triangle-aware interpolation of both channels.\\
\bottomrule
\end{tabularx}
\end{table}

\subsection{Solver workflow}

A call to \texttt{solve\_jet\_rg} performs the following operations:
\begin{enumerate}[leftmargin=1.5em]
  \item Build the descending uniform log-\(b\) lattice and derive one time node
  from each value of \(b_*(b)\).
  \item Evaluate the user boundary function once at every lattice node.
  \item Evaluate \(\alpha_s(\mu_n)\) and precompute the full splitting table
  for every time and shift. For RK2, also build the midpoint table.
  \item March upward in time, activating each node on its boundary point.
  \item March downward by columns to fill the lower triangle.
  \item Store both \(N\times N\) channel histories and return a callable
  \texttt{StepwiseRGSolution}.
\end{enumerate}

\subsection{Minimal use}

The following example matches the current public API. The user-defined
\texttt{b\_sat} controls the saturation of \(b_*(b)\); the separate repository
variable \texttt{b0} remains the QCD constant
\(\bzero=2\ee^{-\gamma_E}\) used internally in \(\mu_*=\bzero/b_*\).

\begin{lstlisting}
include("Numerical RG/numerical_rg.jl")

b_sat = 1.0 # GeV^-1

bstar_func(b) = b / sqrt(1.0 + b^2 / b_sat^2)

boundary_func(; b, bstar, mu_start) = (
    exp(-2.0 * b),        # Jq(b, mu_start)
    0.5 * exp(-3.0 * b), # Jg(b, mu_start)
)

solution = solve_jet_rg(
    n_nodes = 1000,
    b_min = 0.001,
    b_max = 30.0,
    bstar_func = bstar_func,
    boundary_func = boundary_func,
    order = 2,
    method = :rk2,
)
\end{lstlisting}

The boundary function is intentionally generic. It receives the physical
arguments \texttt{b}, \texttt{bstar}, and \texttt{mu\_start}, and must return
the finite pair \texttt{(jq, jg)}. A perturbative boundary calculation and a
nonperturbative model can be inserted later without changing the solver.

\subsection{Inputs}

\begin{table}[H]
\centering
\caption{Public keyword arguments of \texttt{solve\_jet\_rg}.}
\label{tab:solver-inputs}
\begin{tabularx}{\textwidth}{@{}>{\ttfamily}l >{\raggedright\arraybackslash}p{0.18\textwidth} X@{}}
\toprule
\textrm{Argument} & \textrm{Allowed value} & \textrm{Meaning}\\
\midrule
n\_nodes
  & integer \(\ge2\)
  & Number \(N\) of paired \(\ell\) and \(t\) nodes.\\
b\_min
  & positive Float64
  & Smallest physical \(b\), in \(\mathrm{GeV}^{-1}\), represented by the
    lattice.\\
b\_max
  & Float64 \(>\texttt{b\_min}\)
  & Largest represented \(b\), in \(\mathrm{GeV}^{-1}\), beyond which zero
    closure is imposed.\\
bstar\_func
  & function
  & Maps each \(b_i\) to a finite positive \(b_{*,i}\). It must produce
    strictly increasing \(t_i=2\ln(\bzero/b_{*,i})\) in lattice order.\\
boundary\_func
  & function
  & Returns \texttt{(jq, jg)} at \((b_i,b_{*,i},\mu_i)\). The optional default
    returns \texttt{(1.0, 1.0)}.\\
order
  & \(0,1,2\)
  & Highest splitting order: LO, NLO, or NNLO, with all lower orders included.\\
method
  & \texttt{:euler} or \texttt{:rk2}
  & Explicit time integrator. The public default is \texttt{:rk2}.\\
alpha\_s\_Z\_value
  & positive Float64
  & Fixed-\(n_f=5\) reference coupling at \(91.2\,\mathrm{GeV}\); default
    \(0.118\). No threshold matching is applied.\\
alpha\_s\_loops
  & \(1,\ldots,4\)
  & Loop order used by the running coupling; default \(4\).\\
alpha\_s\_max
  & positive Float64
  & Safety ceiling for every sampled \(\alpha_s(\mu)\); default \(1.0\).\\
\bottomrule
\end{tabularx}
\end{table}

There are deliberately no \texttt{mu\_final}, \texttt{t\_start}, or
\texttt{max\_delta\_t} arguments. The complete scale interval and every time
step follow from \(b_{\min}\), \(b_{\max}\), \(N\), and \texttt{bstar\_func}.
This table documents the public \texttt{solve\_jet\_rg} entry point. The
lower-level \texttt{solve\_stepwise\_rg} function has a separate
\texttt{:euler} default and is mainly intended for testing and internal use.

\subsection{Output and interpolation}

The return value is a \texttt{StepwiseRGSolution}. It contains the two lattice
histories, boundary values, lattice metadata, and final active state. For
normal use, treat it as a callable interpolator:

\begin{lstlisting}
both = solution(1.0, 10.0)
jq = both.jq
jg = both.jg

jq_only = solution(1.0, 10.0; component = :q)
jg_only = solution(1.0, 10.0; component = :g)

canonical_mu = mu_start(solution, 1.0)
\end{lstlisting}

Valid queries satisfy
\begin{equation}
  b_{\min}\le b\le b_{\max},
  \qquad
  \ee^{t_1/2}\le\mu\le\ee^{t_N/2}.
  \label{eq:query-domain}
\end{equation}
Queries outside this rectangle raise a \texttt{DomainError}; the interpolator
does not extrapolate. At a stored lattice node it returns the stored value
exactly.

\subsection{Threading and computational cost}

Run Julia with multiple threads:
\begin{lstlisting}
julia --project="Numerical RG" --threads=auto your_script.jl
\end{lstlisting}
The implementation uses conservative parallelism:
\begin{itemize}[leftmargin=1.5em]
  \item kernel-table rows are built in parallel;
  \item independent active-node right-hand sides use at most four Julia tasks;
  \item each individual convolution dot product remains serial and uses
  \texttt{@simd}, avoiding nested threading overhead.
\end{itemize}

Running-coupling settings are stored in repository-level globals before a
kernel table is built. Consequently, solver calls with different coupling
settings must not run concurrently, and the solver should not be nested
inside another static threaded loop.

The kernel tables and solution histories require \(O(N^2)\) memory. Building
the kernel table costs \(O(N^2)\). The present explicit full-plane evolution
performs a growing convolution at many nodes and time slices, giving
\(O(N^3)\) leading work.

The benchmark script performs a warm-up and accepts the node count and BLAS
thread count:
\begin{lstlisting}
cd "Numerical RG"
julia --project=. --threads=20 test/benchmark_solver.jl 1000 1
\end{lstlisting}
Representative recorded timings with 20 Julia threads and one BLAS thread
were:
\begin{equation}
\begin{array}{c|rrrr}
N & 500 & 1000 & 2000 & 4000\\
\hline
\text{wall time [s]} & 0.45 & 3.5 & 32.6 & 286.8
\end{array}
\label{eq:timings}
\end{equation}
The hardware identifier and raw timing logs were not retained, so these values
are comparative development measurements rather than a reproducibility-grade
benchmark. The near-cubic growth is more important than the individual
numbers.

\subsection{Accuracy guidance and practical choices}

The default recommendation is \texttt{method = :rk2}. An exactly solvable
local-kernel diagnostic gives second-order convergence for RK2 and first-order
convergence for Euler on an affine boundary line. It does not test the
nonlocal channel coupling, a nonlinear \(b_*\) curve, or every detail of the
lower-region RK2 construction. The companion test report keeps the
quantitative validation results and explains these limits.

In production use:
\begin{itemize}[leftmargin=1.5em]
  \item choose \(b_{\max}\) so the jet function is already numerically zero at
  the endpoint, then enlarge it until observables are closure-insensitive;
  \item compare at least two values of \(N\);
  \item inspect the minimum canonical scale and ensure
  \(\alpha_s(\mu)\le\texttt{alpha\_s\_max}\);
  \item do not interpret agreement between two coarse grids as an analytic
  validation.
\end{itemize}

\subsection{Tests and reproducibility}

From the repository root, run
\begin{lstlisting}
julia --project="Numerical RG" -e "using Pkg; Pkg.instantiate()"

julia --project="Numerical RG" --threads=auto "Numerical RG/test/runtests.jl"

julia --project="Numerical RG" --threads=auto "Numerical RG/test/analytic_validation.jl"

julia --project="Numerical RG" --threads=auto "Numerical RG/test/momentum_sum_check.jl"
\end{lstlisting}
The regression suite contains 262 asserted checks. The two remaining scripts
are diagnostics: they print convergence and moment residuals but currently do
not fail when a numerical threshold is exceeded. Detailed logic, exact
parameters, and results are kept in \texttt{test\_report.tex}.


\section{Scope and interpretation}

The solver now provides the full-plane object needed for phenomenological
use, but its output remains a numerical approximation to
Eq.~\eqref{eq:bspace-rg}. The main limitations are:
\begin{itemize}[leftmargin=1.5em]
  \item the zero closure at \(b_{\mathrm{grid,max}}\), including its
  incompatibility with a nonzero endpoint value in the plus distributions;
  \item the lattice quadrature in \(u\), whose error can dominate RK2;
  \item the fact that refining \(N\) simultaneously refines the \(\ell\) and
  boundary-derived \(t\) grids;
  \item a fixed, massless \(n_f=5\) running coupling and splitting scheme, with
  no heavy-flavor threshold matching or low-scale freezing;
  \item global running-coupling settings, which make concurrent solver calls
  with different coupling inputs non-reentrant;
  \item the need to keep every sampled \(\mu\) in a perturbative region.
\end{itemize}

The companion \texttt{test\_report.tex} explains the local-kernel convergence
diagnostic, the physical momentum-sum diagnostic, and their limitations. The
two diagnostic scripts currently print measurements but do not enforce
pass/fail thresholds; only \texttt{runtests.jl} is an asserted regression
suite.

\end{document}
